% Удиртгал
% \phantomchapter{Удиртгал} % Зарим нэг зөвлөмж

\phantomsection
\addchaptertocentry{Удиртгал}
\label{Introduction} % Энэ бүлэг рүү ишлэл хийх бол \ref{Introduction} командыг ашигла 
%-------------------------------------------------------------------------------
%   INTRODUCTION
%-------------------------------------------------------------------------------
% Зорилго зорилт, өмнөх бие даалт харах
    \section*{Удиртгал} %for example:
    Орчин үеийн хөдөлмөрийн зах зээлд ажил олгогч нь ажил горилогч болгоны ур чадварыг гараар үнэлэхэд цаг хугацаа их шаардагдаж байна. Мөн ажил горилогчид өөрсдийн ур чадвар тухайн ажлын байранд хэр нийцэж байгааг бодитоор үнэлүүлэх, сул болон давуу талуудаа тодорхой мэдэх боломж хязгаарлагдмал байдаг. Иймд ажил горилогчийн анкет болон ажлын байрны шаардлагыг хиймэл оюунаар харьцуулан үнэлж, тайлбар бүхий дүгнэлт гаргах, цаашдын шатны даалгаврыг автоматаар үүсгэх мөн цаашлаад өөрт тохирох ажлын байрыг чатботоор дамжуулан хурдан шуурхай олж өгөх ухаалаг систем боловсруулах нь энэхүү ажлын үндсэн зорилго болно. Энэхүү систем нь ажил горилогч, ажил олгогч болон хиймэл оюуны гурвалсан харилцан үйлчлэл дээр суурилсан бөгөөд хөдөлмөрийн зах зээл дэх сонгон шалгаруулалтын процессыг илүү бодит, ил тод, хурдан шуурхай болгоход чиглэнэ.
    \subsection*{Зорилго}
    Энэхүү төслийн ажлын зорилго нь ажил олгогч болон ажил горилогчдод зориулсан, ажил горилогчийн ур чадварыг тухайн ажлын байрны шаардлагад тулгуурлан хиймэл оюун ухаанаар үнэлэх веб суурьтай системийг 2026 оны 6 сар хүртэл хөгжүүлнэ. Энэ систем нь хэрэглэгч болон ажлын байрны нийцлийг харуулж, хүний нөөцийн дараагийн шатны ажлуудыг хиймэл оюун ашиглан хөнгөвчилж, автоматжуулах ба ажил горилогч өөрт тохирсон ажлын байрыг цаг тухайд нь, хурдан олж, оновчтой шийдвэр гаргах боломжийг бүрдүүлж өгнө.
    
    \subsection*{Зорилт}
        Уг зорилгод хүрэхийн тулд дараах зорилтуудыг дэвшүүлж байна. Үүнд:
        \begin{enumerate}[nosep]
        \item Ажил горилогч, ажил олгогчийн системийн шаардлагыг тодорхойлох
        \item Ур чадварын үнэлгээний загвар боловсруулах
        \item Анкет болон ажлын байрны шаардлагыг хиймэл оюунаар задлан шинжлэх алгоритм боловсруулах
        \item Ур чадварын тохирцын хувийг тооцоолох механизм боловсруулах
        \item Шалгуур хангаагүй ур чадваруудын тайлбар бүхий дүгнэлт гаргах
        \item Дараагийн шатны даалгаврыг AI-аар үүсгэх болон засварлах модуль боловсруулах
        \item Чатбот ашиглан ажил горилогчид тохиромжтой ажлын байр санал болгох систем хэрэгжүүлэх
        \item Системийн гүйцэтгэл, нарийвчлалыг туршилтаар үнэлэх
        \end{enumerate}

    \subsection*{Судлах үндэслэл}
    Орчин үеийн хөдөлмөрийн зах зээлд байгууллагууд чадварлаг, үр бүтээлтэй ажилтан сонгон шалгаруулах шаардлагатай тулгарч байна. Дижитал шилжилт, технологийн хурдацтай хөгжил, олон улсын хөдөлмөрийн өрсөлдөөн нэмэгдсэнтэй холбоотойгоор ажил горилогчдын ур чадварын шаардлага улам бүр нарийсч, мэргэжлийн болон зөөлөн ур чадвар (soft skills)-ын тэнцвэрт байдал чухал үзүүлэлт болж байна. Гэвч практикт байгууллагууд ажил горилогчийн анкет, CV, ярилцлага зэрэг уламжлалт аргуудад тулгуурлан үнэлгээ хийдэг бөгөөд энэ нь дараах хүндрэлүүдийг дагуулдаг. Үүнд:

        \begin{enumerate}[nosep]
            \item Анкетын мэдээлэл субъектив, өөрийгөө дөвийлгөсөн байх магадлалтай
            \item Олон тооны хүсэлтийг богино хугацаанд шүүхэд хүний нөөцийн зардал өндөр
            \item Ур чадварын бодит тохирцыг тоон болон чанарын хувьд нарийвчлан үнэлэх боломж хязгаарлагдмал
            \item Ярилцлагад үндэслэсэн шийдвэр гаргалт нь хүний субъектив хандлагаас хамаарах
        \end{enumerate}

        Нөгөө талаас ажил горилогчид өөрийн ур чадвар тухайн ажлын байрны шаардлагад хэр нийцэж байгааг урьдчилан мэдэх, сул талаа тодорхойлох, өөрийгөө хөгжүүлэх чиглэлээ тогтоох хэрэгцээ нэмэгдэж байна. Гэвч одоогийн системүүд нь зөвхөн ажлын байрны мэдээлэл нийтлэх, анкет илгээх түвшинд ажиллаж байгаа бөгөөд ур чадварын задлан шинжилгээ, тайлбар бүхий дүгнэлт, автомат зөвлөмж өгөх боломж хомс байна.
        Сүүлийн жилүүдэд хиймэл оюун ухаан, машин сургалтын аргууд нь текстэн өгөгдлийг ойлгож, ангилж, тохирцын оноо тооцоолох, урьдчилсан таамаглал гаргах өндөр чадвартай болсон. Иймд ажил горилогчийн анкет болон ажлын байрны шаардлагыг автоматаар харьцуулж, ур чадварын тохирцыг тооцоолох, тайлбар бүхий дүгнэлт гаргах, цаашдын шатны даалгаврыг автоматаар үүсгэх систем боловсруулах нь онолын болон практикийн хувьд ач холбогдолтой асуудал болж байна. Иймээс энэхүү төгсөлтийн ажлын хүрээнд хиймэл оюунд суурилсан ур чадварын үнэлгээний веб системийг зохион байгуулж, хөдөлмөрийн зах зээл дэх сонгон шалгаруулалтын процессыг илүү оновчтой, ил тод, хурдан шуурхай болгох шаардлага үүссэн болно.

    \subsection*{Судлагдсан байдал}
 
    Сүүлийн жилүүдэд хиймэл оюун ухаанд суурилсан хүний нөөцийн удирдлага, сонгон шалгаруулалтын системүүд эрчимтэй хөгжиж байна. Уламжлалт ярилцлага, анкетын үнэлгээнээс гадна өгөгдөлд суурилсан, алгоритм ашигласан шийдвэр гаргалт нь илүү найдвартай, бодитой байж болохыг дараах судалгаануудыг онцолж байна.
    \begin{itemize}
        \item T. Chamorro-Premuzic, D. Winsborough, R. A. Sherman, R. Hogan нар (2016) судалгаандаа хиймэл оюунд суурилсан үнэлгээний аргууд нь хүний субъектив хандлагаас хамаарах нөлөөллийг багасгаж, илүү тогтвортой, давтагдах боломжтой үр дүн гаргах боломжтойг дурдсан \cite{Chamorro2016}. Тэд өгөгдөлд тулгуурласан үнэлгээ нь ажилтны гүйцэтгэлийг урьдчилан таамаглах чадвартай болохыг онцолсон.
        \item A. K. Upadhyay болон K. Khandelwal (2018) нарын судалгаагаар хиймэл оюун ухаанд суурилсан recruitment системүүд нь анкет шүүх хугацааг бууруулж, хүний нөөцийн зардлыг хэмнэх, олон тооны өргөдлийг богино хугацаанд боловсруулах боломжтойг тогтоосон \cite{Upadhyay2018}. Энэхүү судалгаа нь AI технологи нь зөвхөн автоматжуулалт бус стратегийн түвшний шийдвэр гаргалтад нөлөөлөх боломжтойг харуулсан.
        \item J. Malinowski, T. Keim, O. Wendt, T. Weitzel нар (2006) ажил горилогч болон ажлын байрны шаардлагын хоорондох тохирцыг алгоритмын аргаар тооцоолох хоёр талт зөвлөмжийн (bilateral recommendation) загварыг боловсруулсан бөгөөд энэ нь ур чадварын нийцлийг тооцоолох онолын суурь болж байна \cite{Malinowski2006}.
        \item J. S. Black болон P. van Esch (2020) нар AI-д суурилсан сонгон шалгаруулалтын системийн хэрэглээ, боломж болон эрсдэлийг судалж, байгууллагууд алгоритмын ил тод байдал, ёс зүйн асуудлыг анхаарах шаардлагатайг дурдсан \cite{Black2020}.
    \end{itemize}

    Дээрх судалгаанууд нь хиймэл оюун ухааныг хүний нөөцийн салбарт ашиглах боломжийг онолын болон практикийн түвшинд баталсан боловч ур чадварын тохирцыг тайлбар бүхий задлан шинжилгээ хийх, дараагийн шатны даалгаврыг автоматаар үүсгэх, чатботоор ажлын байр санал болгох цогц системийн шийдэл харьцангуй ховор байна. Иймд энэхүү судалгаа нь AI-д суурилсан ур чадварын үнэлгээний веб системийг цогцоор нь боловсруулж, хөдөлмөрийн зах зээлд практик хэрэглээ болгох зорилготой юм.

    \subsection*{Шинэлэг тал}

    Энэхүү системийн шинэлэг талууд нь дараах байдалтай байна:

    \begin{enumerate}[itemsep=0mm]
    \item Ур чадварын тохирцыг зөвхөн хувиар илэрхийлэхгүй, чадвар тус бүрээр тайлбар бүхий задлан шинжилгээ хийх
    \item Шалгуур хангаагүй ур чадвар бүрийн шалтгааныг хиймэл оюунаар текстэн тайлбар хэлбэрээр гаргах
    \item Ажил олгогчийн шаардлагад тулгуурлан дараагийн шатны даалгаврыг автоматаар үүсгэх боломжтой байх
    \item Ажил горилогчийн хүссэн цалин, байршил, чиглэлд тулгуурлан чатботоор ажлын байр санал болгох
    \item Ажил горилогчийн “ирээдүйн үнэ цэнэ” буюу хөгжих боломжийг харгалзан шийдвэр гаргалтад дэмжлэг үзүүлэх
    \item Ажил горилогч, ажил олгогч, AI гурвалсан харилцан үйлчлэлд суурилсан архитектуртай байх
    \end{enumerate}
        
    \subsection*{Технологийн ач холбогдол}

    Энэхүү систем нь өөрийн машин сургалтын загвар сургах бус, харин бэлэн сургалттай хиймэл оюуны API үйлчилгээг ашиглан текстэн өгөгдөлд боловсруулалт хийх архитектур дээр суурилна. Ийм шийдэл нь өндөр тооцооллын нөөц шаардалгүйгээр орчин үеийн байгалийн хэл боловсруулалтын (NLP) боломжийг практикт нэвтрүүлэх давуу талтай. API-д суурилсан архитектур нь системийн өргөтгөх чадвар, найдвартай ажиллагаа, засвар үйлчилгээний хялбар байдлыг хангадаг. Мөн өгөгдлийн аюулгүй байдал, нууцлалын асуудлыг зохих протокол, шифрлэлтийн технологи ашиглан шийдвэрлэх боломжтой. 

    Иймээс энэхүү шийдэл нь хүний нөөцийн процессыг автоматжуулах, стратегийн шийдвэр гаргалтыг өгөгдөлд тулгууртай болгох, мөн AI технологийг бодит хэрэглээнд нэвтрүүлэх жишиг загвар болох технологийн ач холбогдолтой юм.
    \subsection*{Хамрах хүрээ}
    Энэхүү систем нь веб суурьтай байдлаар дараах үндсэн хэрэглэгчдийг хамарна:
    \begin{itemize}
        \item Ажил горилогч
        \item Ажил олгогч (Хүний нөөц ба байгууллагууд)
    \end{itemize}

    Системийн хүрээнд:
    \begin{itemize}
        \item Ажил горилогч бүртгүүлэх, анкет үүсгэх
        \item Ажлын байрны шалгуур оруулах
        \item Ур чадварын тохирцыг хиймэл оюунаар үнэлэх
        \item Дараагийн шатны даалгавар үүсгэх
        \item Даалгавар хүлээн авч үнэлэх
        \item Чатботоор ажлын байр санал болгох
    \end{itemize}
        зэрэг функцүүд хэрэгжинэ.

    Энэхүү систем нь хөдөлмөрийн зах зээлд оролцогч хувь хүн, хувийн болон төрийн байгууллагууд ашиглах боломжтой.